% CERP Exam Question Bank - Set 3 (Questions 401-600)
% Additional 200 Unique Questions - Not Repeated from Sets 1 or 2
% Compiled from ABA CERP Content Outline
% Compiled on: 2026-06-13

\documentclass[12pt, letterpaper]{article}
\usepackage[utf8]{inputenc}
\usepackage{geometry}
\usepackage{amsmath}
\usepackage{amsfonts}
\usepackage{amssymb}
\usepackage{titlesec}
\usepackage{enumitem}
\usepackage{xcolor}
\usepackage{fancyhdr}
\usepackage{multicol}
\usepackage{longtable}

\geometry{margin=1in}
\setlength{\parindent}{0pt}
\setlength{\parskip}{0.8ex}

\pagestyle{fancy}
\fancyhf{}
\fancyhead[L]{\textbf{CERP Exam Prep - Set 3}}
\fancyhead[R]{\thepage}
\fancyfoot[C]{© 2026 Enterprise Risk Training}

\newcounter{questioncounter}
\setcounter{questioncounter}{401}

\newcommand{\question}{
    \par\noindent
    \textbf{\arabic{questioncounter}.}\quad
    \stepcounter{questioncounter}
}

\newcommand{\answer}[1]{
    \par\noindent
    \textbf{Answer:} #1
    \vspace{0.3cm}
    \par
}

\titleformat{\section}{\large\bfseries}{\thesection}{1em}{}
\titleformat{\subsection}{\normalsize\bfseries}{\thesubsection}{1em}{}

\begin{document}

\begin{center}
    \vspace*{1cm}
    {\huge \textbf{Certified Enterprise Risk Professional (CERP)}} \\[0.5cm]
    {\huge \textbf{Exam Question Bank - Set 3}} \\[1cm]
    {\Large 200 Additional Multiple-Choice Questions (401-600)} \\[1.5cm]
    {\large Based on ABA CERP Exam Content Outline} \\[0.3cm]
    {\large Domains 1-8: Risk Governance \& Risk Management} \\[2cm]
    {\large \textbf{Questions 401 - 600}} \\[0.3cm]
    \vfill
    {\large \today}
\end{center}

\newpage

% ============================================================
% DOMAIN 1: Risk Governance (Questions 401-416)
% ============================================================
\section{RISK GOVERNANCE}

\subsection{Domain 1: Board and Senior Management Oversight (Questions 401-416)}

\question Which document formally defines the risk types and levels the board expects management to maintain?
A) Strategic plan
B) Risk appetite statement
C) Marketing brochure
D) Employee handbook
\answer{B}

\question What is the board's role in risk culture?
A) Delegate entirely to junior staff
B) Model desired behaviors and hold management accountable
C) Ignore culture
D) Focus only on financial metrics
\answer{B}

\question Which is a sign that risk culture is weakening?
A) Increased risk reporting
B) Fear of raising concerns or "shooting the messenger"
C) Higher training attendance
D) More frequent board meetings
\answer{B}

\question What information should management provide to the board regularly?
A) Only good news
B) Risk profile, limit breaches, emerging risks, and control effectiveness
C) Employee vacation schedules
D) Office supply inventory
\answer{B}

\question What is the consequence of ineffective board oversight?
A) Stronger risk management
B) Unchecked risk-taking and potential failure
C) Higher profits guaranteed
D) Lower regulatory scrutiny
\answer{B}

\question How does risk appetite influence strategic decisions?
A) It has no influence
B) It determines which strategies are acceptable based on risk capacity
C) It only affects operational decisions
D) It only affects IT decisions
\answer{B}

\question Which committee typically oversees the whistleblower program?
A) Marketing committee
B) Audit committee
C) Loan committee
D) Facilities committee
\answer{B}

\question What is management's accountability regarding risk appetite?
A) No accountability
B) Operating within the approved risk appetite and reporting breaches
C) Changing appetite without board approval
D) Ignoring appetite
\answer{B}

\question Which behavior indicates a healthy risk culture?
A) Hiding losses
B) Open discussion of near-misses and lessons learned
C) Blaming others for errors
D) Avoiding all risk-taking
\answer{B}

\question What is the role of the Chief Risk Officer (CRO)?
A) Only credit risk
B) Lead the enterprise risk management function independently
C) Replace the CEO
D) Manage marketing
\answer{B}

\question How often should the board assess the effectiveness of the risk management program?
A) Never
B) At least annually or more frequently based on risk
C) Once every 5 years
D) Only after a crisis
\answer{B}

\question What is a "risk committee charter"?
A) A marketing document
B) A document defining the committee's authority, composition, and responsibilities
C) An employee contract
D) A vendor agreement
\answer{B}

\question Which is a key element of risk governance?
A) No defined roles
B) Clear assignment of risk responsibilities across three lines of defense
C) One person does everything
D) No oversight
\answer{B}

\question What is the board's role in risk appetite review?
A) Approve without review
B) Critically review and approve changes to risk appetite
C) Delegate entirely to the CRO
D) Ignore appetite
\answer{B}

\question Which is a leading practice for board risk reporting?
A) Dense, 500-page reports
B) Concise dashboards with key metrics, trends, and exception highlights
C) No reports
D) Verbal updates only
\answer{B}

\question What is the relationship between risk capacity and risk appetite?
A) Same thing
B) Risk appetite must be less than or equal to risk capacity
C) Risk capacity is always smaller
D) No relationship
\answer{B}

% ============================================================
% DOMAIN 2: Policies, Procedures and Limits (Questions 417-440)
% ============================================================

\subsection{Domain 2: Policies, Procedures and Limits (Questions 417-440)}

\question What is a "trading limit" for market risk?
A) A suggestion
B) A maximum notional or loss limit for trading activities
C) An employee limit
D) A marketing limit
\answer{B}

\question Who has authority to approve a policy exception above a certain threshold?
A) Any employee
B) Senior management or board depending on materiality
C) External auditor
D) Customer
\answer{B}

\question What is a "stop-loss limit"?
A) A limit on profits
B) A limit that triggers automatic position reduction after a loss level
C) A limit on employees
D) A marketing limit
\answer{B}

\question Which is a typical policy for conflict of interest?
A) No policy
B) Disclosure and recusal requirements
C) Encouraging conflicts
D) Ignoring conflicts
\answer{B}

\question What is a "limit breach"?
A) Following the limit
B) Exceeding an approved risk limit
C) A new policy
D) An approved exception
\answer{B}

\question What is the purpose of a "delegation of authority" matrix?
A) To concentrate authority
B) To specify who can approve what type and size of transaction or exception
C) To eliminate authority
D) To confuse staff
\answer{B}

\question Which is a regulatory expectation for third-party risk policies?
A) No oversight
B) Due diligence, monitoring, and contingency plans for critical vendors
C) Ignore vendor risk
D) No written agreements
\answer{B}

\question What is a "policy inventory"?
A) A list of office supplies
B) A centralized list of all policies with review dates and owners
C) A marketing list
D) An employee list
\answer{B}

\question Which is an example of a risk limit for concentration risk?
A) No limit
B) Maximum 25\% of capital to a single industry sector
C) Unlimited to one borrower
D) No diversification requirement
\answer{B}

\question What is a "procedure" in risk management?
A) A high-level principle
B) Step-by-step instructions to execute a policy
C) A marketing slogan
D) An employee benefit
\answer{B}

\question How should policy violations be tracked?
A) Verbally only
B) In a centralized issue log with remediation plans
C) Ignored if small
D) Only tracked by the violator
\answer{B}

\question What is a "self-assessment" of policy compliance?
A) No assessment
B) Business line's own evaluation of adherence to policies
C) External audit only
D) Regulatory review only
\answer{B}

\question Which is a limit for operational risk?
A) Loan-to-value ratio
B) Maximum acceptable operational loss per event
C) Interest rate duration
D) Credit spread
\answer{B}

\question What is the consequence of operating without risk limits?
A) Lower risk
B) Potential for uncontrolled losses
C) Guaranteed compliance
D) Higher efficiency
\answer{B}

\question Who approves a new risk policy?
A) The policy owner alone
B) Appropriate governance body (e.g., Risk Committee, Board)
C) External consultant
D) Any employee
\answer{B}

\question What is a "risk limit framework"?
A) A single limit
B) The structured set of all risk limits across the organization
C) A marketing plan
D) An HR policy
\answer{B}

\question Which is a limit for country risk?
A) No limit
B) Maximum exposure to a single country based on risk rating
C) Unlimited exposure
D) Only domestic exposure
\answer{B}

\question What is a "breach remediation plan"?
A) No plan
B) Documented actions to address a limit breach and prevent recurrence
C) A marketing plan
D) An employee recognition plan
\answer{B}

\question Which is a typical policy for insider trading?
A) Encouraging insider trading
B) Prohibition and blackout periods
C) No policy
D) Verbal guidance only
\answer{B}

\question What is a "limit monitoring report"?
A) A marketing report
B) A report showing current usage versus approved limits
C) An HR report
D) A facilities report
\answer{B}

\question Which is a limit for liquidity risk?
A) Loan-to-value
B) Minimum stock of High-Quality Liquid Assets (HQLA)
C) Employee count
D) Number of branches
\answer{B}

\question What is a "policy exception"?
A) Following policy exactly
B) Approved deviation from a policy requirement
C) A new policy
D) A deleted policy
\answer{B}

\question Which is a regulatory expectation for model risk policy?
A) No validation
B) Model inventory, validation, and governance
C) Use any model without review
D) No documentation
\answer{B}

\question What is the purpose of a "risk limit escalation policy"?
A) To hide breaches
B) To define when and to whom limit breaches are reported
C) To approve all breaches automatically
D) To ignore breaches
\answer{B}

% ============================================================
% DOMAIN 3: MIS and Data Governance (Questions 441-455)
% ============================================================

\subsection{Domain 3: Management Information Systems (Questions 441-455)}

\question What is a "data quality rule"?
A) No rules
B) A defined standard for data accuracy, completeness, and timeliness
C) A marketing slogan
D) An employee rule
\answer{B}

\question Which is a risk of manual data entry?
A) Perfect accuracy
B) Higher error rates and lack of audit trail
C) Faster processing
D) Lower cost
\answer{B}

\question What is "data reconciliation"?
A) Deleting data
B) Comparing data across systems to ensure consistency
C) Creating duplicate data
D) Ignoring data differences
\answer{B}

\question What is a "critical data element"?
A) Unimportant data
B) Data essential for risk measurement and reporting
C) Marketing data
D) HR data only
\answer{B}

\question What is the purpose of a data governance policy?
A) To allow unrestricted data changes
B) To define roles, standards, and procedures for managing data
C) To delete all data
D) To ignore data quality
\answer{B}

\question Which is a control for system access?
A) No passwords
B) Role-based access control and periodic access reviews
C) Everyone has admin access
D) No logging of access
\answer{B}

\question What is "data lineage" used for?
A) Deleting data
B) Tracing data from source to report to understand transformations
C) Creating new data
D) Ignoring data sources
\answer{B}

\question Which is a risk of spreadsheet-based risk reporting?
A) Low error risk
B) High risk of formula errors, version control issues, and manual manipulation
C) Perfect audit trail
D) Automatic validation
\answer{B}

\question What is a "data validation rule"?
A) A suggestion
B) An automated check to ensure data meets quality standards
C) A marketing rule
D) An employee rule
\answer{B}

\question What is the role of the Data Governance Council?
A) No role
B) Oversee data strategy, quality, and issue resolution
C) Enter data manually
D) Delete data
\answer{B}

\question Which is a KRI for data quality?
A) Number of employees
B) Percentage of data fields with null values or errors
C) Loan balance
D) Branch count
\answer{B}

\question What is a "data dictionary" used for?
A) Storing actual data
B) Defining metadata: field names, types, definitions, and allowed values
C) Deleting data
D) Marketing
\answer{B}

\question Which is a best practice for risk data aggregation?
A) Siloed databases
B) Standardized data taxonomies and centralized data warehouse
C) No data standards
D) Manual aggregation only
\answer{B}

\question What is the consequence of poor data quality for risk reporting?
A) Perfect decisions
B) Misinformed risk decisions and potential regulatory action
C) Lower costs
D) Faster reporting
\answer{B}

\question What is a "data quality dashboard"?
A) No dashboard
B) A visual display of data quality metrics and issues
C) An HR dashboard
D) A marketing dashboard
\answer{B}

% ============================================================
% DOMAIN 4: Control Framework (Questions 456-470)
% ============================================================

\subsection{Domain 4: Control Framework (Questions 456-470)}

\question What is a "manual control"?
A) An automated control
B) A control performed by a person rather than a system
C) No control
D) A marketing control
\answer{B}

\question What is an "automated control"?
A) A manual control
B) A control performed by a system (e.g., system limit check)
C) No control
D) A verbal control
\answer{B}

\question Which is a control for financial reporting under SOX?
A) No controls
B) Reconciliation, authorization, and segregation of duties
C) Unlimited access
D) No documentation
\answer{B}

\question What is a "control deficiency"?
A) A control that works perfectly
B) A weakness in control design or operation
C) A strong control
D) A redundant control
\answer{B}

\question What is a "material weakness" in internal controls?
A) A minor issue
B) A deficiency that results in a reasonable possibility of material misstatement
C) An immaterial issue
D) A control that works
\answer{B}

\question Which is a control for model risk?
A) No model governance
B) Independent model validation and ongoing monitoring
C) Using unvalidated models
D) No documentation
\answer{B}

\question What is a "control owner"?
A) No one
B) The person responsible for maintaining and evidencing a control
C) External auditor
D) Customer
\answer{B}

\question Which is a detective control for fraud?
A) Pre-approval of all transactions
B) Surprise audits and reconciliations
C) Segregation of duties
D) Physical security
\answer{B}

\question What is "control testing frequency" based on?
A) Random choice
B) Risk assessment and control criticality
C) Never test
D) Test once only
\answer{B}

\question Which is a control for business continuity risk?
A) No plan
B) Documented business continuity plan and regular testing
C) Rely on luck
D) No backup systems
\answer{B}

\question What is a "control matrix"?
A) A marketing tool
B) A document mapping risks to controls and testing status
C) An HR form
D) A facilities plan
\answer{B}

\question Which is a preventive control for data entry?
A) After-the-fact review
B) Input validation and field limits in the system
C) Monthly reconciliation
D) Investigation of errors
\answer{B}

\question What is "control automation" intended to achieve?
A) Increase human error
B) Improve consistency, speed, and reduce errors
C) Eliminate all controls
D) Reduce compliance
\answer{B}

\question Which is a control for vendor risk?
A) No due diligence
B) Third-party risk assessment and ongoing monitoring
C) No contract
D) Ignore vendor performance
\answer{B}

\question What is the purpose of a "control self-assessment" (CSA)?
A) To eliminate controls
B) To empower process owners to evaluate their own controls
C) To outsource all control testing
D) To ignore control effectiveness
\answer{B}

% ============================================================
% DOMAIN 5: Risk Identification (Questions 471-495)
% ============================================================

\subsection{Domain 5: Risk Identification (Questions 471-495)}

\question Which is an environmental risk?
A) Loan default
B) Climate change impact on operations or collateral
C) System failure
D) Employee fraud
\answer{B}

\question What is a "risk workshop" used for?
A) Marketing
B) Collaborative identification of risks with stakeholders
C) HR training
D) Facilities planning
\answer{B}

\question Which is a social risk?
A) Interest rate change
B) Changing consumer preferences or labor unrest
C) Credit default
D) System outage
\answer{B}

\question What is a "risk statement"?
A) A marketing slogan
B) A clear description of a risk including cause, event, and impact
C) An employee goal
D) A budget item
\answer{B}

\question Which is a technology risk?
A) Loan default
B) Legacy system failure or cyber attack
C) Interest rate change
D) Deposit withdrawal
\answer{B}

\question What is "horizon scanning"?
A) Looking at the sun
B) Systematic identification of emerging risks and trends
C) Ignoring future risks
D) Only looking at past risks
\answer{B}

\question Which is a risk from regulatory change?
A) No impact
B) New compliance requirements and potential business model changes
C) Guaranteed profit
D) Lower costs
\answer{B}

\question What is a "risk taxonomy" used for?
A) Categorizing animals
B) Standardizing risk language across the organization
C) Marketing segmentation
D) Employee classification
\answer{B}

\question Which is an example of a "black swan" event?
A) Expected loan default
B) A highly improbable event with massive impact
C) Routine operational loss
D) Daily market fluctuation
\answer{B}

\question What is a "risk register"?
A) A marketing document
B) A living document listing identified risks, assessments, and action plans
C) An HR file
D) A facilities log
\answer{B}

\question Which is a political risk?
A) Loan default
B) Expropriation or change in government policy
C) System downtime
D) Employee error
\answer{B}

\question What is "risk identification" the first step of?
A) Marketing
B) The risk management process
C) HR process
D) Facilities management
\answer{B}

\question Which is a source of risk from mergers and acquisitions?
A) Guaranteed success
B) Integration risk, culture clash, and due diligence gaps
C) No risk
D) Only financial gain
\answer{B}

\question What is a "risk appetite" input to risk identification?
A) Ignored
B) Risks outside appetite need identification and response
C) Only risks within appetite matter
D) No relationship
\answer{B}

\question Which is a key output of risk identification?
A) Marketing plan
B) List of risks with preliminary categorization
C) Final financial statements
D) Employee schedule
\answer{B}

\question What is a "risk driver"?
A) A control
B) A factor that increases or decreases risk (e.g., complexity, change)
C) A risk response
D) A risk limit
\answer{B}

\question Which is an emerging risk in banking?
A) Traditional loan default
B) Crypto-asset risks and AI model risks
C) Historical interest rate risk
D) Mature credit risk only
\answer{B}

\question What is the value of "external data" in risk identification?
A) No value
B) Provides insights on industry trends, competitor failures, and regulatory focus
C) Only internal data matters
D) External data is always wrong
\answer{B}

\question Which is a risk from remote work?
A) No risk
B) Cybersecurity and operational resilience challenges
C) Guaranteed productivity
D) Lower costs only
\answer{B}

\question What is a "risk trigger event"?
A) A routine transaction
B) A specific occurrence that activates a pre-defined risk response
C) A marketing campaign
D) An employee meeting
\answer{B}

\question Which is a risk from rapid technological change?
A) No risk
B) Obsolescence and skill gaps
C) Guaranteed efficiency
D) Lower costs only
\answer{B}

\question What is "brainstorming" in risk identification?
A) A medical procedure
B) A facilitated group session to generate risk ideas
C) A solo activity
D) A marketing technique
\answer{B}

\question Which is a risk from climate change for a bank?
A) No risk
B) Physical risk to collateral and transition risk to borrowers
C) Only reputational risk
D) Only operational risk
\answer{B}

\question What is a "risk profile"?
A) A marketing document
B) A summary of an entity's risk exposures and risk management maturity
C) An employee evaluation
D) A budget
\answer{B}

\question Which is a source of risk from outsourcing IT?
A) No risk
B) Vendor lock-in, data security, and service disruption
C) Guaranteed cost savings
D) Higher internal efficiency
\answer{B}

% ============================================================
% DOMAIN 6: Risk Measurement (Questions 496-520)
% ============================================================

\subsection{Domain 6: Risk Measurement and Evaluation (Questions 496-520)}

\question What is "risk quantification"?
A) Describing risk in words only
B) Assigning numerical values to likelihood and impact
C) Ignoring numbers
D) Only qualitative assessment
\answer{B}

\question Which is a measure of credit risk?
A) Duration
B) Probability of Default (PD) and Loss Given Default (LGD)
C) Value at Risk (VaR)
D) System uptime
\answer{B}

\question What is "expected shortfall" (ES)?
A) The same as VaR
B) The average loss beyond the VaR confidence level
C) The minimum loss
D) The median loss
\answer{B}

\question Which is a measure of operational risk under Basel?
A) Loan-to-value
B) Loss Distribution Approach (LDA) or Standardized Approach
C) Duration
D) Credit spread
\answer{B}

\question What is "risk-weighted assets" (RWA)?
A) Total assets
B) Assets adjusted for risk to determine capital requirements
C) Book value of assets
D) Market value of assets
\answer{B}

\question Which is a key metric for liquidity risk?
A) Loan-to-deposit ratio
B) Liquidity Coverage Ratio (LCR) and Net Stable Funding Ratio (NSFR)
C) Return on Equity
D) Efficiency ratio
\answer{B}

\question What is "sensitivity analysis" used for?
A) No analysis
B) Assessing how changes in inputs affect outputs
C) Only qualitative analysis
D) Ignoring variables
\answer{B}

\question Which is a measure of concentration risk?
A) HHI (Herfindahl-Hirschman Index)
B) Total assets
C) Number of employees
D) Branches count
\answer{A}

\question What is "scenario analysis"?
A) Predicting the future exactly
B) Evaluating portfolio performance under specific hypothetical conditions
C) Historical backtesting only
D) No analysis
\answer{B}

\question Which is a measure of interest rate risk?
A) Loan growth
B) Economic Value of Equity (EVE) sensitivity
C) Deposit balance
D) Employee turnover
\answer{B}

\question What is "credit migration risk"?
A) No risk
B) Risk of credit rating downgrades
C) Risk of default only
D) Risk of prepayment
\answer{B}

\question Which is a measure of model risk?
A) Number of models
B) Model validation status and performance monitoring results
C) Employee count
D) Branch count
\answer{B}

\question What is "basis risk"?
A) No risk
B) Risk that imperfect hedging leads to residual exposure
C) Credit default risk
D) Operational risk
\answer{B}

\question Which is a key metric for profitability risk?
A) Loan default rate
B) Net Interest Margin (NIM) volatility
C) System downtime
D) Employee satisfaction
\answer{B}

\question What is "risk aggregation" used for?
A) Ignoring correlations
B) Understanding total enterprise risk across business lines
C) Only single risk focus
D) Marketing
\answer{B}

\question Which is a measure of reputational risk?
A) Stock price only
B) Media sentiment analysis and stakeholder surveys
C) Loan default rate
D) Operational loss amount
\answer{B}

\question What is "value at risk" (VaR)?
A) Guaranteed loss
B) Maximum expected loss over a given horizon at a confidence level
C) Minimum loss
D) Average loss
\answer{B}

\question Which is a limitation of VaR?
A) Too simple
B) Does not measure tail risk beyond confidence level
C) Always accurate
D) No assumptions needed
\answer{B}

\question What is "stressed VaR"?
A) Normal VaR
B) VaR calculated using stressed market parameters
C) Minimum VaR
D) Ignoring stress
\answer{B}

\question Which is a measure of operational resilience?
A) Loan growth
B) Recovery Time Objectives (RTO) and Recovery Point Objectives (RPO)
C) Deposit balance
D) Employee count
\answer{B}

\question What is "risk factor sensitivity"?
A) No sensitivity
B) How much a portfolio value changes per unit change in a risk factor
C) Only qualitative
D) Ignoring risk factors
\answer{B}

\question Which is a measure of counterparty credit risk?
A) Loan-to-value
B) Potential Future Exposure (PFE) and Current Exposure (CE)
C) System uptime
D) Employee turnover
\answer{B}

\question What is "wrong-way risk"?
A) No risk
B) Correlation between exposure and counterparty credit quality that increases risk
C) Right-way risk
D) Operational risk
\answer{B}

\question Which is a measure of liquidity gap?
A) Loan-to-deposit ratio
B) Cumulative mismatch of cash inflows and outflows over time buckets
C) Capital ratio
D) ROE
\answer{B}

\question What is "risk appetite utilization"?
A) Ignoring appetite
B) Current risk levels compared to approved appetite limits
C) Only qualitative
D) Marketing metric
\answer{B}

% ============================================================
% DOMAIN 7: Risk Responses (Questions 521-545)
% ============================================================

\subsection{Domain 7: Risk Responses (Questions 521-545)}

\question When is "risk mitigation" the appropriate response?
A) When risk is already within appetite
B) When risk exceeds appetite and can be cost-effectively reduced
C) Never
D) Only for reputational risk
\answer{B}

\question What is "risk transfer" through hedging?
A) Ignoring risk
B) Using derivatives to offset market risk exposure
C) Accepting all risk
D) Avoiding all activity
\answer{B}

\question What is a "root cause"?
A) The symptom
B) The underlying reason a problem occurred
C) The final impact
D) The loss amount
\answer{B}

\question Which is an example of risk mitigation for cyber risk?
A) No controls
B) Firewalls, intrusion detection, and employee training
C) Buying insurance only
D) Ignoring cyber risk
\answer{B}

\question What is "risk avoidance" for credit risk?
A) Lending to everyone
B) Declining to lend to a high-risk borrower
C) Accepting the loan
D) Transferring via credit default swap
\answer{B}

\question What is an "issue remediation plan"?
A) No plan
B) Documented actions to fix an identified deficiency
C) A marketing plan
D) An HR plan
\answer{B}

\question Which is an example of risk acceptance?
A) Adding controls
B) Self-insuring and booking reserves for expected losses
C) Avoiding the activity
D) Hedging
\answer{B}

\question What is "risk response effectiveness"?
A) No measurement
B) Whether the response reduced risk to within appetite
C) Cost only
D) Speed only
\answer{B}

\question Which is an example of risk transfer via insurance?
A) Self-insurance
B) Purchasing a policy to cover operational losses
C) Avoiding all risk
D) Accepting all risk
\answer{B}

\question What is a "corrective action"?
A) No action
B) Action taken to fix a problem and prevent recurrence
C) A preventive action
D) A detective action
\answer{B}

\question What is the purpose of "root cause analysis" in issue management?
A) To blame someone
B) To fix the underlying problem, not just symptoms
C) To ignore the issue
D) To delay action
\answer{B}

\question Which is an example of risk mitigation for compliance risk?
A) Ignoring regulations
B) Implementing compliance monitoring and training programs
C) Accepting all compliance violations
D) Avoiding all business
\answer{B}

\question What is "residual risk rating"?
A) Inherent risk rating
B) Risk rating after applying risk responses
C) No rating
D) Control rating only
\answer{B}

\question Which is an example of a risk response for a low-priority risk?
A) Extensive mitigation
B) Monitor only or accept within appetite
C) Immediate avoidance
D) Expensive transfer
\answer{B}

\question What is the goal of issue validation?
A) To assume issue is fixed
B) To independently confirm remediation was effective
C) To ignore closure
D) To close without evidence
\answer{B}

\question Which is an example of risk mitigation for fraud risk?
A) No controls
B) Segregation of duties and mandatory vacations
C) Accepting all fraud
D) Avoiding all transactions
\answer{B}

\question What is a "risk response plan"?
A) No plan
B) A documented strategy for addressing identified risks
C) A marketing plan
D) An HR plan
\answer{B}

\question Which is an example of risk avoidance for regulatory risk?
A) Violating regulations
B) Ceasing a non-compliant product or activity
C) Accepting fines
D) Transferring via insurance
\answer{B}

\question What is the role of the first line in risk response?
A) Audit the response
B) Implement and own risk responses
C) Set risk appetite
D) Validate responses
\answer{B}

\question What is "issue aging"?
A) Age of employees
B) How long an issue has remained open without remediation
C) Age of customers
D) Age of policies
\answer{B}

\question Which is an example of risk transfer via contract?
A) No contract
B) Indemnification clauses shifting liability to vendor
C) Accepting all liability
D) Avoiding contracts
\answer{B}

\question What is the purpose of a "risk response cost-benefit analysis"?
A) Ignore costs
B) Ensure response cost does not exceed risk reduction benefit
C) Always choose most expensive response
D) Never analyze
\answer{B}

\question Which is an example of risk mitigation for third-party risk?
A) No due diligence
B) Regular vendor assessments and contractual SLAs
C) Accepting all vendor failures
D) Avoiding all vendors
\answer{B}

\question What is "residual risk tolerance"?
A) No tolerance
B) The acceptable level of risk after responses are applied
C) Inherent risk level
D) Control effectiveness
\answer{B}

\question Which is an example of risk acceptance with monitoring?
A) No monitoring
B) Accepting low risks but tracking them for change
C) Ignoring changes
D) Only accepting without monitoring
\answer{B}

% ============================================================
% DOMAIN 8: Risk Monitoring (Questions 546-570)
% ============================================================

\subsection{Domain 8: Risk Monitoring (Questions 546-570)}

\question What is a "key performance indicator" (KPI) in risk context?
A) A risk metric
B) A metric measuring performance of risk processes (e.g., issue closure time)
C) A marketing metric
D) An HR metric
\answer{B}

\question What is a "risk threshold"?
A) A suggestion
B) A specific value that triggers escalation or action
C) A final loss
D) A marketing target
\answer{B}

\question Which is a KRI for model performance?
A) Number of employees
B) Backtesting exceptions or model drift
C) Branch count
D) Loan balance
\answer{B}

\question What is "ongoing monitoring" vs. periodic monitoring?
A) Same thing
B) Continuous vs. at set intervals
C) Periodic is better
D) Ongoing is only for IT
\answer{B}

\question Which is a KRI for vendor risk?
A) Loan growth
B) Number of vendor audit findings overdue
C) Deposit balance
D) Employee count
\answer{B}

\question What is a "risk dashboard" intended to provide?
A) Too much detail
B) At-a-glance view of risk posture and exceptions
C) No information
D) Only historical data
\answer{B}

\question Which is a KRI for business continuity risk?
A) System uptime only
B) Time since last BCP test or number of failed test scenarios
C) Employee satisfaction
D) Marketing spend
\answer{B}

\question What is "limit utilization"?
A) Ignoring limits
B) Current exposure divided by approved limit (percentage)
C) The limit amount
D) The breach amount
\answer{B}

\question Which is a KRI for conduct risk?
A) Loan growth
B) Number of customer complaints or regulatory referrals
C) Deposit balance
D) Capital ratio
\answer{B}

\question What is the purpose of a "risk committee report"?
A) To hide risks
B) To summarize top risks, limit breaches, and emerging issues for committee review
C) To market products
D) To evaluate employees
\answer{B}

\question Which is a KRI for capital adequacy?
A) Loan-to-deposit ratio
B) CET1 ratio and capital conservation buffer
C) Employee turnover
D) Customer satisfaction
\answer{B}

\question What is "risk escalation"?
A) Ignoring issues
B) Raising a risk or issue to a higher level of management for decision
C) Deleting the risk
D) Accepting the risk without action
\answer{B}

\question Which is a KRI for money laundering risk?
A) Loan growth
B) Number of suspicious activity reports (SARs) filed
C) Deposit balance
D) Branch count
\answer{B}

\question What is a "risk appetite breach"?
A) Operating within appetite
B) Exceeding a defined risk appetite limit
C) A new policy
D) An approved exception
\answer{B}

\question Which is a KRI for stress testing?
A) Number of employees
B) Number of scenarios not meeting capital adequacy under stress
C) Loan balance
D) Marketing spend
\answer{B}

\question What is "control monitoring frequency" based on?
A) Random selection
B) Risk assessment and control criticality
C) Never monitor
D) Only monitor after loss
\answer{B}

\question Which is a KRI for audit issue remediation?
A) Number of new loans
B) Number of overdue audit remediation actions
C) Deposit growth
D) Employee count
\answer{B}

\question What is a "risk trend analysis"?
A) Looking at one point
B) Analyzing risk metrics over time to identify deterioration
C) Ignoring trends
D) Only current snapshot
\answer{B}

\question Which is a KRI for climate risk?
A) Loan growth
B) Percentage of portfolio in climate-vulnerable sectors or geographies
C) Employee count
D) Branch count
\answer{B}

\question What is the purpose of "exception reporting"?
A) To report all transactions
B) To focus attention on deviations from policy or limits
C) To hide deviations
D) To report only normal activity
\answer{B}

\question Which is a KRI for internal fraud?
A) Loan growth
B) Number of policy overrides without business justification
C) Deposit balance
D) Capital ratio
\answer{B}

\question What is a "risk and control monitoring plan"?
A) No plan
B) Documented schedule and methods for monitoring key risks and controls
C) A marketing plan
D) An HR plan
\answer{B}

\question Which is a KRI for third-party concentration?
A) Number of employees
B) Percentage of critical activities dependent on a single vendor
C) Loan balance
D) Branch count
\answer{B}

\question What is the role of the second line in risk monitoring?
A) Own all monitoring
B) Oversee and challenge first-line monitoring
C) No role
D) Only monitor annually
\answer{B}

\question Which is a KRI for change management risk?
A) Number of loans
B) Number of unauthorized or failed changes to production systems
C) Deposit balance
D) Employee satisfaction
\answer{B}

% ============================================================
% FINAL QUESTIONS 571-600 (Mixed Domains)
% ============================================================

\subsection{Comprehensive Review (Questions 571-600)}

\question What is the difference between "risk appetite" and "risk limit"?
A) Same thing
B) Appetite is broad strategic; limits are specific quantitative boundaries
C) Limits are broader
D) Appetite is quantitative only
\answer{B}

\question Which is a key principle of the three lines of defense?
A) No separation
B) Independence of the third line (audit)
C) First line does audit
D) Second line has no role
\answer{B}

\question What is "operational risk" under Basel?
A) Credit risk only
B) Risk of loss from inadequate or failed internal processes, people, systems
C) Market risk only
D) Strategic risk only
\answer{B}

\question Which is a key output of an RCSA?
A) Marketing plan
B) Residual risk ratings and action plans for risks above appetite
C) Employee schedule
D) Budget
\answer{B}

\question What is the purpose of a "risk control library"?
A) No purpose
B) Centralized repository of control definitions and attributes
C) Marketing tool
D) Employee directory
\answer{B}

\question Which is a regulatory expectation for capital planning?
A) No planning
B) Stress testing and capital adequacy assessment based on risk profile
C) Ignore capital
D) Only regulatory capital minimum
\answer{B}

\question What is "risk transparency"?
A) Hiding risks
B) Clear visibility into risk exposures, assumptions, and decisions
C) No reporting
D) Only positive risks shared
\answer{B}

\question Which is a key element of risk culture?
A) No accountability
B) Psychological safety to raise concerns
C) Punishment for all errors
D) No training
\answer{B}

\question What is "model validation" independence?
A) Same team builds and validates
B) Validator independent from model development
C) No validation
D) Vendor validation only
\answer{B}

\question Which is a key metric for interest rate risk in the banking book (IRRBB)?
A) Loan growth
B) EVE sensitivity and NII sensitivity
C) Deposit balance
D) Employee count
\answer{B}

\question What is the purpose of a "risk report" to the Board?
A) To overwhelm with data
B) To inform decision-making and oversight
C) To hide negative information
D) To market products
\answer{B}

\question Which is a key component of the COSO ERM framework?
A) Only control activities
B) Governance and culture, strategy and objective-setting, performance, review and revision, information and communication
C) Only monitoring
D) Only risk assessment
\answer{B}

\question What is "risk assessment" the foundation for?
A) Marketing
B) Risk response and control design
C) HR decisions
D) Facilities management
\answer{B}

\question Which is a key output of scenario analysis?
A) Exact predictions
B) Identification of vulnerabilities and capital needs under stress
C) No insights
D) Only qualitative
\answer{B}

\question What is the purpose of a "risk policy"?
A) To restrict all activity
B) To establish principles, roles, and requirements for managing risk
C) To eliminate risk
D) To market products
\answer{B}

\question Which is a key metric for operational risk?
A) Loan default rate
B) Operational loss frequency and severity
C) Interest rate duration
D) Credit spread
\answer{B}

\question What is "risk reporting to regulators"?
A) Optional
B) Mandatory submission of risk data as required by law
C) Never required
D) Only upon request
\answer{B}

\question Which is a key element of risk governance?
A) No board involvement
B) Clear roles, responsibilities, and accountability
C) One person makes all decisions
D) No documentation
\answer{B}

\question What is the purpose of a "risk and control self-assessment"?
A) To eliminate controls
B) To evaluate risk and control effectiveness by process owners
C) To outsource risk management
D) To ignore risks
\answer{B}

\question Which is a key metric for compliance risk?
A) Loan growth
B) Number of regulatory findings and remediation status
C) Deposit balance
D) Employee count
\answer{B}

\question What is "risk reporting frequency" based on?
A) Random
B) Risk level and audience (e.g., daily for trading, quarterly for Board)
C) Always annual
D) Never report
\answer{B}

\question Which is a key principle of internal control?
A) No segregation of duties
B) Accountability and documentation
C) Unlimited access
D) No monitoring
\answer{B}

\question What is the purpose of a "risk appetite dashboard"?
A) To hide breaches
B) To visually show risk levels versus appetite limits
C) To market products
D) To evaluate employees
\answer{B}

\question Which is a key output of issue management?
A) No action
B) Remediated issues with validated closure
C) Unresolved issues
D) Ignored findings
\answer{B}

\question What is the purpose of "risk training"?
A) To waste time
B) To ensure employees understand risks and their roles
C) To eliminate risk
D) To market products
\answer{B}

\question Which is a key metric for strategic risk?
A) Loan default rate
B) Variance from strategic plan milestones
C) System uptime
D) Employee satisfaction
\answer{B}

\question What is the role of the second line in policy management?
A) Own all policies
B) Oversee and challenge policy adherence and exceptions
C) No role
D) Only implement
\answer{B}

\question Which is a key component of a risk register?
A) Employee birthdays
B) Risk description, owner, assessment, and action plan
C) Marketing slogans
D) Office locations
\answer{B}

\question What is the purpose of "risk appetite metrics"?
A) No purpose
B) To quantify and monitor adherence to risk appetite
C) To ignore limits
D) To confuse staff
\answer{B}

\question Which is a key element of a control framework?
A) No monitoring
B) Control activities and segregation of duties
C) Unlimited exceptions
D) No documentation
\answer{B}

% ============================================================
% ANSWER KEY (Questions 401-600)
% ============================================================
\newpage
\section*{Answer Key - Questions 401-600}
\begin{multicols}{3}
\small
\begin{enumerate}[start=401]
\item B
\item B
\item B
\item B
\item B
\item B
\item B
\item B
\item B
\item B
\item B
\item B
\item B
\item B
\item B
\item B
\item B
\item B
\item B
\item B
\item B
\item B
\item B
\item B
\item B
\item B
\item B
\item B
\item B
\item B
\item B
\item B
\item B
\item B
\item B
\item B
\item B
\item B
\item B
\item B
\item B
\item B
\item B
\item B
\item B
\item B
\item B
\item B
\item B
\item B
\item B
\item B
\item B
\item B
\item B
\item B
\item B
\item B
\item B
\item B
\item B
\item B
\item B
\item B
\item B
\item B
\item B
\item B
\item B
\item B
\item B
\item B
\item B
\item B
\item B
\item B
\item B
\item B
\item B
\item B
\item B
\item B
\item B
\item B
\item B
\item B
\item B
\item B
\item B
\item B
\item B
\item B
\item B
\item B
\item B
\item B
\item B
\item B
\item B
\item B
\item B
\item B
\item B
\item B
\item B
\item B
\item B
\item B
\item B
\item B
\item B
\item B
\item B
\item B
\item B
\item B
\item B
\item B
\item B
\item B
\item B
\item B
\item B
\item B
\item B
\item B
\item B
\item B
\item B
\item B
\item B
\item B
\item B
\item B
\item B
\item B
\item B
\item B
\item B
\item B
\item B
\item B
\item B
\item B
\item B
\item B
\item B
\item B
\item B
\item B
\item B
\item B
\item B
\item B
\item B
\item B
\item B
\item B
\item B
\item B
\item B
\item B
\item B
\item B
\item B
\item B
\item B
\item B
\item B
\item B
\item B
\item B
\item B
\item B
\item B
\item B
\item B
\item B
\item B
\item B
\end{enumerate}
\end{multicols}

\end{document}